\documentclass[11pt]{article}
\usepackage[a4paper,margin=2.5cm]{geometry}
\usepackage[T1]{fontenc}
\usepackage[utf8]{inputenc}
\usepackage{amsmath,amssymb}
\usepackage{booktabs,array}
\usepackage{xcolor}
\usepackage{listings}
\usepackage[numbers,sort&compress]{natbib}
\usepackage[colorlinks=true,linkcolor=blue!50!black,citecolor=blue!50!black,urlcolor=blue!50!black]{hyperref}
\lstset{basicstyle=\ttfamily\small,columns=fullflexible,frame=single,framerule=0.3pt,
  rulecolor=\color{black!30},xleftmargin=4pt,xrightmargin=4pt}

\title{Luz polarizada y fase geométrica en unidades qg:\\
parámetros de Stokes, matrices de Mueller y teorema de Stokes}
\author{Vicente Humberto Monteverde\\
\small Universidad del Museo Social Argentino (UMSA), Argentina\\
\small ORCID 0000-0001-8884-4811}
\date{\small Borrador del 3 de octubre de 2026}

\renewcommand{\tablename}{Tabla}
\renewcommand{\refname}{Referencias}
\renewcommand{\abstractname}{Resumen}

\begin{document}
\sloppy
\maketitle

\begin{abstract}
El proyecto qang describe los qubits con cantidades de medición, los valores qg
$qg_P=\langle P\rangle$~\cite{monteverde2026qang}. Esta nota breve expresa dos resultados clásicos en esas unidades,
implementados y verificados en la librería \texttt{qang} (\texttt{pip install qang}, versión 0.6.0). Primero,
un estado de polarización de la luz es un qubit: los parámetros de Stokes normalizados son exactamente los
valores qg, el grado de polarización es la longitud del vector qg, y una matriz de Mueller sin pérdidas es la
regla de compuertas qg $qg'_P=qg_{U^\dagger PU}$. La ley de Malus queda $I=I_0(1+qg_Z\cos2\theta+qg_X\sin2\theta)/2$
para cualquier estado de entrada. Segundo, la fase geométrica de un lazo cerrado en la esfera de Bloch
(Poincaré) es menos la mitad del ángulo sólido encerrado, que es el teorema de Stokes para la curvatura de
Berry. Para un lazo a $qg_Z$ constante la fase es $(qg_Z-1)/2$ de vuelta: la fija $qg_Z$ sola. Nada de esto es
física nueva. El aporte es la traducción a las unidades de qang y un conjunto de 26 verificaciones numéricas,
con la fase calculada de tres maneras independientes.
\end{abstract}

\section{Los parámetros de Stokes son valores qg}

Con $|H\rangle=|0\rangle$ y $|V\rangle=|1\rangle$, los parámetros de Stokes~\cite{stokes1852} de un estado de
polarización $\rho$ son $S_\mu=I_0\,\mathrm{Tr}(\rho\,\sigma_\mu)$ con $\sigma=(I,Z,X,Y)$. Entonces
\begin{equation}
qg_Z=\frac{S_1}{S_0}=\frac{I_H-I_V}{S_0},\qquad
qg_X=\frac{S_2}{S_0}=\frac{I_D-I_A}{S_0},\qquad
qg_Y=\frac{S_3}{S_0}=\frac{I_R-I_L}{S_0},
\end{equation}
con $|D\rangle,|A\rangle=(|H\rangle\pm|V\rangle)/\sqrt2$ y $|R\rangle,|L\rangle=(|H\rangle\pm i|V\rangle)/\sqrt2$.
Los textos difieren en cuál estado circular llaman derecho. La librería fija $R$ como arriba, así que $S_3>0$
significa $qg_Y>0$, y ofrece la convención opuesta como opción. La esfera de Poincaré es la esfera de Bloch. La
luz parcialmente polarizada es un estado mixto: su grado de polarización $P=\sqrt{S_1^2+S_2^2+S_3^2}/S_0$ es la
longitud del vector qg, y su pureza es $\mathrm{Tr}\rho^2=(1+P^2)/2$. Un despolarizador $\mathrm{diag}(1,d,d,d)$
es el canal despolarizante.

\paragraph{Matrices de Mueller.} Para una matriz de Jones $J$, $M_{\mu\nu}=\tfrac12\mathrm{Tr}(\sigma_\mu J\sigma_\nu
J^\dagger)$. En el orden $(I,Z,X,Y)$ es la matriz de transferencia de Pauli. Para un elemento sin pérdidas
(lámina de onda, rotador), $J$ es unitaria, y la acción de $M$ sobre el vector de Stokes es la regla de
compuertas qg de \texttt{qang.formulation}. Una lámina de media onda a ángulo $\theta$ lleva polarización
lineal a ángulo $\alpha$ a $2\theta-\alpha$: una rotación de $4\theta-2\alpha$ en el plano $(qg_Z,qg_X)$.

\paragraph{Ley de Malus en unidades qg.} Detrás de un polarizador lineal ideal a ángulo $\theta$, la intensidad
transmitida es
\begin{equation}
I(\theta)=\frac{I_0}{2}\left(1+qg_Z\cos2\theta+qg_X\sin2\theta\right),
\end{equation}
que se reduce a $I_0\cos^2\theta$ para luz horizontal. Como $I(\theta)$ depende de la entrada solo a través de
$qg_Z$ y $qg_X$, dos posiciones del polarizador leen dos valores qg y una lámina de cuarto de onda da el tercero.

\paragraph{Desde conteos.} \texttt{qg\_from\_counts} convierte los conteos del analizador
$(n_H,n_V,n_D,n_A,n_R,n_L)$ en los tres valores qg y $P$, con intervalos de \texttt{qang.statistics.qg\_estimate}
(bayesiano, Wilson o método delta). La polarimetría de fotones individuales es entonces el mismo problema de
estimación que la tomografía de qubits en unidades qg.

\section{Fase geométrica y teorema de Stokes}

Un qubit llevado por un lazo cerrado $C$ en la esfera de Bloch adquiere una fase
geométrica~\cite{pancharatnam1956,berry1984}
\begin{equation}
\gamma(C)=-\tfrac12\,\Omega(C),
\end{equation}
donde $\Omega$ es el ángulo sólido encerrado por $C$. Es el teorema de Stokes. La conexión de Berry
$A=i\langle\psi|\nabla\psi\rangle$ tiene curvatura $F=\nabla\times A$, un campo de monopolo de intensidad $1/2$ en
el centro de la esfera. Su integral de línea alrededor de $C$ es igual a su flujo a través del casquete
encerrado, $\Omega/2$. Para un lazo a $qg_Z=\cos\vartheta$ constante alrededor del eje $Z$,
$\Omega=2\pi(1-\cos\vartheta)=2\pi(1-qg_Z)$, así que en la unidad de fase qg$_\Phi$ (fracciones de vuelta)
\begin{equation}
\phi=\frac{\gamma}{2\pi}=\frac{qg_Z-1}{2}\pmod 1 .
\end{equation}
La fase geométrica de un lazo cónico la fija $qg_Z$ sola. Para luz polarizada, la misma fase es la de
Pancharatnam en la esfera de Poincaré. Por ejemplo, el lazo $H\to D\to R\to H$ encierra un octante
($\Omega=\pi/2$) y da $\gamma=-\pi/4$.

\paragraph{Tres cálculos independientes.} \texttt{qang.geometric} calcula la fase de tres maneras:
\begin{enumerate}
\item el producto de solapamientos invariante de gauge $\gamma=-\arg\prod_k\langle\psi_k|\psi_{k+1}\rangle$ sobre
  un lazo discreto de estados~\cite{samuel1988};
\item el ángulo sólido con signo del lazo, como abanico de triángulos esféricos~\cite{vanoosterom1983};
\item el flujo numérico de la curvatura a través del casquete.
\end{enumerate}
Para un lazo discreto, la fase de solapamientos es exactamente $-\tfrac12$ del ángulo sólido del polígono
geodésico. El flujo y la integral de línea coinciden, como exige el teorema de Stokes.

\section{Verificaciones}

\begin{table}[h]
\centering\small
\begin{tabular}{p{0.62\textwidth}l}
\toprule
verificación & resultado\\
\midrule
los seis estados base $H,V,D,A,R,L$ dan $\pm1$ en un eje qg & exacto\\
ida y vuelta Stokes $\leftrightarrow$ qg, contra \texttt{qang.formulation.qg\_values} (20 estados al azar) & $10^{-12}$\\
Mueller sin pérdidas = regla de compuertas qg (4 elementos $\times$ 5 estados) & $10^{-12}$\\
ley de Malus para $H$, $D$, $R$ y un estado elíptico, 7 ángulos & $10^{-12}$\\
la lámina de media onda lleva la polarización lineal a $2\theta-\alpha$ & $10^{-12}$\\
grado de polarización y pureza tras un despolarizador (10 estados) & $10^{-12}$\\
fase de solapamientos $=-\tfrac12$ ángulo sólido geodésico (4 alturas) & $10^{-10}$\\
fase del cono $=(qg_Z-1)/2$ vueltas (4 alturas, 4000 puntos) & $10^{-5}$\\
flujo de curvatura $=\pi(1-qg_Z)=$ integral de línea (Stokes, 3 casquetes) & $10^{-5}$\\
invariancia de gauge; cambio de signo al invertir el lazo & $10^{-12}$\\
octante de Pancharatnam $H\to D\to R$: $\Omega=\pi/2$, $\gamma=-\pi/4$ & exacto\\
\bottomrule
\end{tabular}
\caption{Verificaciones numéricas en \texttt{tests/test\_polarization.py} (12 pruebas) y
\texttt{tests/test\_geometric.py} (14 pruebas).}
\end{table}

\section{Alcance}

Los resultados son óptica y mecánica cuántica estándar. ¿Para qué escribirlos en unidades qg? Las herramientas
de qang para qubits se aplican sin cambios a los qubits fotónicos de polarización. Entre ellas están la
estimación con intervalos desde conteos, la regla de compuertas como tabla de transformaciones qg y la unidad
de fase qg$_\Phi$. Las ecuaciones de Maxwell en sí no forman parte de esta nota: describen campos clásicos,
mientras que qang mide cantidades de qubits.

\section*{Código}

\begin{lstlisting}[language=Python]
from qang import polarization as P, geometric as G
P.stokes_to_qg([1.0, 0.6, 0.0, 0.8])     # {'X': 0.0, 'Y': 0.8, 'Z': 0.6}
P.malus_intensity({"Z": 1.0}, 0.3)        # cos(0.3)**2
G.cone_phase_turns(0.5)                   # 0.75 de vuelta = -pi/2 rad
\end{lstlisting}

\bibliographystyle{unsrtnat}
\bibliography{../refs}

\end{document}
